Legislation is pushing automotive manufacturers to adopt increasingly complex autonomous driving systems,
on a trajectory that will ultimately require full, real-time spatial understanding of the surrounding environment.
As a result, detection pipelines have grown substantially in complexity and compute demand,
while vehicles remain subject to tight on-board power and thermal budgets.
These trends suggest that scaling conventional dense DNNs on GPU-class accelerators is increasingly unsustainable,
motivating ultra-efficient architectures such as spiking neural networks for tasks in resource-constrained environments.

We investigate the viability of spiking networks for automotive-grade perception, and measure the cost of the spiking paradigm against a continuous equivalent.
We propose DetectorNX: a directly-trained spiking detector (DetectorNX-G3-SNN) designed around activation sparsity to minimise energy consumption while preserving detection accuracy, 
paired with a structurally equivalent continuous baseline (DetectorNX-G3-CNN) at identical parameter count,
trained on the ETraM event-camera dataset to isolate the spiking mechanism as the sole variable.
To the best of our knowledge, this is the first directly-trained SNN developed and evaluated on ETraM.

DetectorNX-G3-SNN achieves $94.9\%$ mIoU parity and $91.3\%$ $mAP_{50}$ parity against the continuous baseline,
with $100\%$ recall at IoU $> 0.5$.
Per-frame theoretical energy consumption drops by $94.33\%$, scaling linearly with scene activity ($R^2 = 0.7514$).
The SNN maintains $259.80$ FPS on GPU hardware ($8.6\times$ the automotive real-time threshold) with superior confidence calibration in low-confidence regions.

This project establishes directly-trained SNNs as a viable route to energy-efficient automotive-grade vehicle detection, at a cost of only $\sim5\%$ \gls{miou} precision under strict architectural \gls{snn}--\gls{cnn} parity. 
The $1$-bit quantisation penalty underlying this gap acts as a scale-independent global ceiling, invalidating the initial scale-dependent hypothesis. 
Superior low-confidence calibration further positions the SNN as an implicit uncertainty sensor, safety-relevant for autonomous driving. 
DetectorNX thereby serves as an ETraM baseline, and proposes areas for future work to close the gap with published broader-class spiking detectors.